\clearpage
\fancyhead[L]{\small STERNBEDECKUNG.}
\hypertarget{native-453-sternbedeckung}{}
\begin{center}
{\large STERNBEDECKUNG.}

\rule{0.75\linewidth}{0.4pt}

Astronomische Nachrichten. Band III. Nr. 50. Seite 22. 1824 Februar.

\rule{0.75\linewidth}{0.4pt}
\end{center}

\begin{center}
Bedeckung der Plejaden am 29. August 1820.
\end{center}

Dieselbe Bedeckung ist in Göttingen, Bremen, Altona, Hamburg, Berlin, Wien und Modena beobachtet.

\begin{center}
Göttingen (in einem Briefe von \textsc{Gauss} [an \textsc{Bessel}] freundlicherweise mitgeteilt).

\begingroup\small\setlength{\tabcolsep}{4pt}
\begin{tabular}{c l c c c | c c l}
$p$ & Austritt & $9^{\mathrm u}$ & $46^{\mathrm m}$ & $42^{\mathrm s}49$ & $10^{\mathrm u}$ & $5^{\mathrm m}$ & $26^{\mathrm s}78 - 1.582x$ \\
$\eta$ & --- & $9$ & $51$ & $37.37$ & $10$ & $8$ & $22.89 - 1.265x$ \\
$f$ & Eintritt & $9$ & $40$ & $4.77$ & $10$ & $46$ & $56.01 + 0.219x$ \\
$h$ & --- & $9$ & $44$ & $11.60$ & $10$ & $49$ & $18.09 + 0.795x$ \\
\end{tabular}
\endgroup
\end{center}

[wobei $x$ die Korrektur der Breitendifferenz zwischen Mond und Stern bezeichnet.]

